(1) 标量维里定理 \(2K+W=0\) 给 \(W=-2K\)。因此系统总能
\[
E=K+W=-K=-\frac32Nk_BT,\qquad W=-3Nk_BT.
\]
总能量为负，并随 T 上升而进一步降低。

(2) 沿 N 固定且持续维里平衡的状态序列定义有效热容
\[
C_{\rm eff}=\frac{dE}{dT}=-\frac32Nk_B<0.
\]
注意单粒子动能的正定热容仍是 \(3Nk_B/2\)，负号来自同时变化的引力势能。

(3) 当辐射带走能量，\(\dot E=-L\)，代入 \(E=-3Nk_BT/2\) 得
\[
\dot T=\frac{2L}{3Nk_B}>0.
\]
系统向外散热却升温，并不违反能量守恒：引力收缩释放的势能超过辐射损失，其中一部分转为粒子的随机热运动动能。

(4) 若 \(W=-\alpha GM^2/R\)，由维里定理
\[
K=\frac{\alpha GM^2}{2R},\quad
T=\frac{\alpha GM^2}{3Nk_BR},\quad
E=-\frac{\alpha GM^2}{2R}.
\]
半径减小时温度按 R 的倒数增长。由 Ri 收缩到 Rf<Ri 时，
\[
E_{\rm rad}=E_i-E_f=
\frac{\alpha GM^2}{2}\left(\frac1{R_f}-\frac1{R_i}\right)>0.
\]
这些结论依赖于近似维里平衡、可忽略外表面压力和结构系数恒定的前提。
